\section{Methodological spine}
\label{sec:method}

This section is a placeholder for the full method. The final paper should present the method as an evidence-bounded replay pipeline rather than as an unconstrained LLM reasoning chain. The recommended subsection structure is:

\subsection{Task definition}
Define Evidence-Bounded Driver Process-Model Replay. State the inputs, outputs, admissible claims, and forbidden claims.

\subsection{Evidence representation}
Describe reported, derived, not-reported, narrative, and counterfactual evidence. State explicitly that not reported is not evidence of absence.

\subsection{STPA-HF driver process-model representation}
Define CPS, CPB, OPS, and OPB. Explain why these variables mediate the transition from scene facts to driver-action hypotheses.

\subsection{Process-model update and action selection}
Describe process-model formation/update sources, observed update issues, abductive update hypotheses, blocked update claims, other factors, and candidate driver actions.

\subsection{UCA pathway generation and ranking}
Explain that UCA pathways are derived forward from process model and action candidates, while reported outcomes are used only as compatibility constraints. Describe deterministic evidence gates and LLM judge ranking.

\subsection{Replay package output}
Describe the final tabletop replay package, including replay questions, blocked claims, ranked pathways, and minimum HMI/logging requirement candidates.

\subsection{Claim boundary}
State that the method does not infer true driver psychology, prove real HMI causality, reconstruct a unique accident cause, assign legal responsibility, or treat pathway scores as causal probabilities.
